\scanpage{103}\mspage{- 166 -}\hfill IX

qui n'est autre que 12.5 dans le cas de $\underline X^\lambda$ sur $\underline S^\lambda$. Or $\underline S$ est un trait de \uline{caractéristique nulle}. Cela achève donc la réduction annoncée.

\uline{Remarques} 12.9. On notera que cette réduction nous réduit en même temps au cas où $X_{\unsure{o|0}}$ n'a qu'un seul point singulier, qui est celui qui était considéré classiquement dans la formule de Picard-Lefschetz. Notons d'ailleurs que dans ce cas, $X_o$ a une ou deux composantes irréductibles, ces cas correspondants respectivement à $\operatorname{rang} M=1$ et $\operatorname{rang} M=0$ (ce dernier cas étant celui où la jacobienne de la fibre générique a bonne réduction).

12.10. La démonstration donnée dans 12.8 nous fournit un résultat plus général que 12.5, valable sans supposer que $X$ soit régulier, en introduisant les entiers $d_\lambda\geq1$ de (12.8.12) : on trouve alors que les accouplements de monodormie se déduisent tous de l'accouplement
\[
 u=\sum_{\lambda\in L}d_\lambda\varphi_{x_\lambda}\otimes\varphi_{x_\lambda}:M\unsure{\otimes|\hat\otimes}M\longrightarrow\mathbb Z\quad.
 \tag*{(12.10.1)}
\]

13. \uline{Liens avec la théorie transcendante} : \uline{cas analytique complexe}.

Pour les notions générales relatives aux espaces analytiques, nous renvoyons à [2].

13.1. \uline{Rappels sur les groupes analytiques sur} $S$. Soit $S$ un espace analytique, $A$ un espace analytique au-dessus de $S$, muni d'une structure de groupe commutatif au-dessus de $S$ (i.e. $A$ est un groupe commutatif de la catégorie des espaces analytiques au-dessus de $S$). On suppose
% archive sheet number: 102
% stamp: numbered label 484; 484 also in pencil above the lower strip.
